% ======================================================================
\section{Orientation: what you are looking at}
% ======================================================================

Before the instructions, a map. CONCLAVE replaces one component of a conventional
AI agent --- the language model at the centre --- with a committee of small
machine-learning models. Everything else an agent needs is still there.

\begin{center}
\begin{tikzpicture}[
  font=\sffamily\footnotesize,
  box/.style={draw=conclaveblue!55, fill=conclavebg, rounded corners=2pt,
              minimum height=8.5mm, inner xsep=5pt, align=center},
  core/.style={draw=conclavelight, fill=conclavelight!12, line width=1pt,
               rounded corners=2pt, minimum height=11mm, inner xsep=6pt,
               align=center},
  gate/.style={draw=conclaveamber, fill=conclaveamber!10, rounded corners=2pt,
               minimum height=8.5mm, inner xsep=5pt, align=center},
  ar/.style={-{Stealth[length=2mm]}, conclaveblue!70, line width=0.8pt}]

\node[box] (obs)  {Observation\\{\scriptsize sensors + text}};
\node[box, right=7mm of obs] (feat) {Featurise\\{\scriptsize hashing + scaling}};
\node[core, right=7mm of feat] (brain) {\textbf{Ensemble brain}\\{\scriptsize 3--4 small models}};
\node[box, right=7mm of brain] (mem) {Episodic\\prior};
\node[gate, right=7mm of mem] (gt) {Uncertainty\\gate};
\node[box, right=7mm of gt] (act) {Act / abstain\\escalate};

\draw[ar] (obs)--(feat); \draw[ar] (feat)--(brain); \draw[ar] (brain)--(mem);
\draw[ar] (mem)--(gt);   \draw[ar] (gt)--(act);

\node[draw=conclaveteal!60, fill=conclaveteal!8, rounded corners=2pt,
      below=8mm of brain, minimum height=7mm, inner xsep=5pt,
      align=center, font=\sffamily\scriptsize] (con)
      {Hard constraints \& planner};
\draw[ar, conclaveteal] (con) -- (gt);
\draw[ar, conclaveteal, dashed] (brain) -- (con);

\begin{scope}[on background layer]
  \node[fit=(obs)(act), draw=conclaverule, line width=0.8pt, rounded corners=3pt,
        inner sep=7pt] {};
\end{scope}
\node[anchor=west, font=\sffamily\scriptsize\bfseries, color=conclavegrey]
      at ([yshift=6.5mm,xshift=-4mm]obs.north west) {ONE TICK OF THE AGENT LOOP};
\end{tikzpicture}
\end{center}

\subsection{The seven objects you will actually touch}

\begin{center}
\renewcommand{\arraystretch}{1.28}
\begin{tabular}{@{}>{\raggedright\arraybackslash}p{3.4cm}
                 >{\raggedright\arraybackslash}p{7.0cm}
                 >{\raggedright\arraybackslash}p{5.1cm}@{}}
\toprule
\rowcolor{conclavebg}
{\sffamily\bfseries\color{conclaveblue}Object} &
{\sffamily\bfseries\color{conclaveblue}What it is} &
{\sffamily\bfseries\color{conclaveblue}Lives in} \\
\midrule
\code{Observation}   & One reading of the world: numeric features plus optional text & \file{conclave/core/types.py} \\
\code{FusedFeaturizer} & Turns observations into vectors under a fixed memory budget & \file{conclave/brain/featurize.py} \\
\code{EnsembleBrain} & The committee. Replaces the LLM & \file{conclave/brain/ensemble.py} \\
\code{UncertaintyGate} & Converts disagreement into act / gather / abstain / escalate & \file{conclave/brain/uncertainty.py} \\
\code{ToolRegistry}  & Your capabilities, with typed schemas & \file{conclave/core/registry.py} \\
\code{Agent}         & The loop that ties it together & \file{conclave/core/agent.py} \\
\code{ToolRouter}    & Picks a tool by ranking, not classifying & \file{conclave/brain/router.py} \\
\bottomrule
\end{tabular}
\end{center}

\begin{keybox}[The one idea worth carrying through this document]
A single model gives you a prediction and a confidence number. A \emph{committee}
gives you something a single model cannot: the \textbf{spread of opinion}
between members. CONCLAVE treats that spread as a first-class control signal ---
it is the difference between ``this is a hard call'' and ``I have no idea'',
and those two situations warrant different behaviour. Section~\ref{sec:math}
derives it properly.
\end{keybox}

\subsection{If you have twenty minutes}

\begin{enumerate}[leftmargin=*,itemsep=3pt]
  \item \S\ref{sec:windows} --- set up the virtual environment, run \code{pytest}.
        Expect \textbf{14 passed}.
  \item Run \code{python demo.py episode} and watch the \code{net} column.
        Long stretches say \code{DOWN} and the agent keeps working. That is the
        entire thesis of the project, visible in one column.
  \item Run \code{python demo.py uncertainty} to see the committee split apart
        as you break its sensors.
\end{enumerate}

Everything else in this guide is depth on those three things.
